\Qhead{3.3}{Модели случайных процессов}{Модели случайных процессов: гауссовский процесс, узкополосный стационарный шум, импульсные случайные процессы, дробовой шум.}
\Subfile{3.3.1}{sections/03_statistics/q03/s01}
\Subfile{3.3.2}{sections/03_statistics/q03/s02}

\recap{Гауссовость задаёт все многовременные распределения через первые два момента. Узкополосная огибающая удобна для амплитуды и фазы; пуассоновские импульсы описывают дискретность независимых событий.}
\board{Две гауссовы квадратуры и Рэлеевская огибающая; сумма импульсов; $G_i=2q\overline i$.}
\oral{Гарантирует ли отсутствие корреляции независимость отсчётов произвольного шума?}
\bait{Можно оставить различие среднего заряда и среднего квадрата: случайное усиление импульсов повышает шум даже при прежнем среднем токе.}
\sources{Биленко и соавт.~\cite{bilenko}, §§4.1.3, 4.2.2; Сёрккя и Солин~\cite{sarkka}, определения гауссовских процессов. Распределение огибающей и импульсная ковариация выведены здесь.}
